% !PRINTED 416
\weberpage{416}
%<<CONT
keinen größeren gemeinsamen Teiler haben als 2. Haben $x$ und $y$ den
größten gemeinschaftlichen Teiler 2, so sind $b$ und $c$ gerade und
$a$, $b$, $c$, $\partial$ sind dann und nur dann relativ prim, wenn $a$
und $\partial$, also auch $n = a\,\partial - b\,c$ \gesperrt{ungerade} sind.

Die Gleichung (8) hat, wenn wir von dem interesselosen Fall $n=1$
absehen, keine diesen Bedingungen genügende Lösung, in der $x=0$ ist.
Ändern wir aber die Vorzeichen von $x$, $y$ zugleich, so gehen $a$,
$b$, $c$, $\partial$ nach (10) in $-a$, $-b$, $-c$, $-\partial$ über, % [sic] source reference (10); (7) is mathematically expected.
und die Gleichung (1) ändert sich nicht. Die Transformation
$\begin{pmatrix} a, & b\\ c, & \partial \end{pmatrix}$ bleibt also
ungeändert.

Wir nennen jetzt der Kürze wegen eine Lösung der Gleichung (8) eine
eigentliche Lösung und $n = \frac{1}{4}(y^2-D\,x^2)$ eine
\gesperrt{eigentliche Darstellung} von $n$, wenn $x$ und $y$ den
folgenden Bedingungen genügt: % [sic] source prints singular genügt with plural subject.

\gesperrt{$x$ ist positiv.}

\gesperrt{Der größte gemeinschaftliche Teiler von $x$ und $y$ ist 1 oder 2.}

\gesperrt{Ist er $=2$, so ist $n$ ungerade.}

Wenn nun $\omega$ einer gegebenen Gleichung (3) genügt, so ist durch
jedes System $\begin{pmatrix} a, & b\\ c, & \partial \end{pmatrix}$
eine eigentliche Lösung von (8) eindeutig bestimmt. Denn nach (5) ist
$x$ der größte gemeinschaftliche Teiler von $b$, $c$, $a-\partial$, und
durch (6) ist $y$ bestimmt, und andererseits ist aus einer eigentlichen
Lösung von (8) durch (7) das System $a$, $b$, $c$, $\partial$
vollständig bestimmt. Es ist aber noch die Frage zu entscheiden, ob
zwei \gesperrt{verschiedene} eigentliche Lösungen $x$, $y$ und $x'$, $y'$ zu
äquivalentem System $\begin{pmatrix} a, & b\\ c, & \partial
\end{pmatrix}$, $\begin{pmatrix} a', & b'\\ c', & \partial'
\end{pmatrix}$ führen können. Unter äquivalentem System sind hier nach
\S\,28, \S\,53 zwei solche zu verstehen, bei denen
\weq{9}{\begin{pmatrix} a', & b'\\ c', & \partial' \end{pmatrix} = \begin{pmatrix} \alpha, & \beta\\ \gamma, & \delta \end{pmatrix}\begin{pmatrix} a. & b\\ c, & \partial \end{pmatrix}} % [sic] source prints a full stop after a in the rightmost matrix.
ist, wenn $\begin{pmatrix} \alpha, & \beta\\ \gamma, & \delta
\end{pmatrix}$ eine \gesperrt{lineare} Transformation ist, also
\weq{10}{\alpha\,\delta - \beta\,\gamma = 1.}
Wenn wir beide Seiten der Gleichung (9) von rechts mit
$\begin{pmatrix} \partial, & -b\\ -c, & a \end{pmatrix}$
zusammensetzen, so ergibt sich:
%CONT>>
